% \documentclass[a4paper,12pt]{article}

% % --- PACOTES FUNDAMENTAIS ---
% \usepackage[utf8]{inputenc}
% \usepackage[T1]{fontenc}
% \usepackage[brazilian]{babel}
% \usepackage{amsmath}
% \usepackage{amssymb}
% \usepackage{geometry}
% \usepackage{graphicx}
% \usepackage{setspace}
% \usepackage{fancyhdr}
% \usepackage{xcolor} % Para cores personalizadas

% % --- CONFIGURAÇÃO DE LINKS (HYPERREF) ---
% \usepackage{hyperref} 

% % Definição de azul escuro profissional (opcional)
% \definecolor{darkblue}{rgb}{0.0, 0.0, 0.55}

% \hypersetup{
%     colorlinks=true,        % Ativa cores nos links
%     linkcolor=black,        % Cor dos links internos (Sumário, referências). Preto para elegância.
%     filecolor=magenta,      
%     urlcolor=blue,          % Cor de URLs externas (GitHub, LinkedIn)
%     citecolor=black,        % Cor de citações bibliográficas
%     pdftitle={Otimização de Eletropostos (MILP) - Proposta Preliminar},
%     pdfpagemode=FullScreen,
%     linktoc=all             % Torna a linha inteira do sumário clicável
% }

% % --- CONFIGURAÇÃO DE MARGENS ---
% \geometry{left=2.5cm,right=2.5cm,top=2.5cm,bottom=2.5cm}

% \begin{document}

% % --- CAPA ---
% \begin{titlepage}
%     \centering
    
%     % Cabeçalho Institucional
%     {\scshape \Large \textbf{UNIVERSIDADE ESTADUAL DE CAMPINAS (UNICAMP)}}\\[0.5cm]
%     {\large Faculdade de Engenharia Elétrica e de Computação (FEEC)}\\[2.0cm]

%     % Título do Projeto
%     \rule{\linewidth}{0.5mm} \\[0.4cm]
%     { \huge \bfseries Proposta preliminar: \\[0.3cm] Otimização da localização de infraestrutura de recarga para veículos elétricos }\\[0.2cm]
%     { \Large Abordagem via Programação Linear Inteira Mista (MILP) com mitigação de canibalização espacial }\\[0.4cm]
%     \rule{\linewidth}{0.5mm} \\[2.5cm]
    
%     % Informações dos Pesquisadores
%     \begin{minipage}{0.8\textwidth}
%         \centering
%         \textbf{Pesquisadores:}\\[0.4cm]
%         {\Large Carlos Julian Muñoz Quiroga}\\[0.2em]
%         \textit{Doutorando em Engenharia Elétrica e de Computação}\\[0.8cm]
        
%         {\Large Núbia Sidene Almeida das Virgens}\\[0.2em]
%         \textit{Revisão e co-desenvolvimento da etapa de otimização}\\
%     \end{minipage}
    
%     \vspace{2cm}
    
%     % Links (Opcional - Pode remover se o repositório for privado)
%     \textbf{Repositório do protótipo e dashboard:}\\
%     \href{https://github.com/CAMUZQUI/designacao-eletropostos-milp}{designacao-eletropostos-milp}
    
%     \vfill
    
%     % Rodapé
%     {\large Campinas, SP}\\
%     {\large \today}
    
% \end{titlepage}

% % --- SUMÁRIO ---
% \newpage
% \thispagestyle{empty} 

% \tableofcontents 

% \newpage

% % --- CONTEÚDO MODULAR ---

% \section{Introdução e contexto do problema}
% \label{sec:introducao}
% Este guia orienta a configuração e utilização da ferramenta de otimização de forma direta e autônoma.

\subsection{Pré-requisitos do sistema}

Para executar esta ferramenta, seu ambiente deve atender a dois requisitos fundamentais:

\subsubsection{Python 3.10}
Verifique se o Python está instalado abrindo o \textit{Prompt de Comando} (cmd) e digitando:

\begin{verbatim}
py --list
\end{verbatim}

\textbf{Como interpretar:}
\begin{itemize}
    \item Se aparecer \texttt{ - 3.10.X} (ou superior), o Python está pronto.
    \item Se der erro, baixe o instalador em \texttt{python.org} e marque a opção "Add to PATH" durante a instalação.
\end{itemize}

\subsubsection{IBM ILOG CPLEX Optimization Studio}
O modelo utiliza a biblioteca \texttt{docplex}, que requer o motor de otimização da IBM instalado no sistema.

\begin{itemize}
    \item É obrigatório ter uma instalação do CPLEX com licença acadêmica ou comercial.
    \item Não utilize a versão "Community Edition", ela possui limites restritos de variáveis e restrições que são insuficientes para a complexidade deste case, impedindo a resolução.
\end{itemize}


\subsection{Configuração do ambiente}

Este passo é necessário apenas na primeira vez que você for usar a ferramenta. Localize e execute o arquivo \texttt{criar\_ambiente.bat} na pasta raiz do projeto.
\begin{itemize}
    \item Uma janela do terminal se abrirá e instalará automaticamente as dependências.
    \item Aguarde até ver a mensagem "AMBIENTE CRIADO COM SUCESSO" e feche a janela.
    \item Se a janela fechar instantaneamente sem mensagens de sucesso, verifique sua instalação do Python.
\end{itemize}

\subsection{Execução da ferramenta}

Para utilizar o otimizador no dia a dia:

\begin{enumerate}
    \item Execute o arquivo \texttt{executar\_solver.bat} na pasta raiz.
    \item Uma janela de seleção de arquivos abrirá automaticamente. Selecione a planilha Excel com os dados de entrada.
    \item Acompanhe o progresso no terminal. O solver processará os dados e buscará a melhor solução.
    \item Ao finalizar, a ferramenta abrirá automaticamente a planilha de resultados gerada na pasta \texttt{resultados}.
\end{enumerate}

\vspace{0.3cm}

\textbf{Controle de Execução (Importante):}

O solver foi configurado para buscar a solução ótima sem um limite de tempo pré-fixado.
\begin{itemize}
    \item \textbf{Interrupção manual:} Você pode interromper o processo a qualquer momento pressionando o atalho \texttt{Ctrl + C} na janela do terminal.
    \item \textbf{Salvamento automático:} Ao interromper o processo, o programa salvará automaticamente a melhor solução encontrada até aquele instante, garantindo que nenhum progresso válido seja perdido.
\end{itemize}


\subsection{Estrutura de arquivos}

A organização das pastas do projeto segue o padrão abaixo. Mantenha essa estrutura para garantir o funcionamento dos scripts automáticos.

\begin{verbatim}
desafio-otimizacao-suzano/
|
+-- criar_ambiente.bat          (Configuracao inicial)
+-- executar_solver.bat         (Execucao diaria)
+-- requirements.txt            (Lista de dependencias)
|
+-- code/                       (Codigos-fonte Python)
|   +-- modelo.py
|   +-- gestor_dados.py
|   +-- exportar_resultados.py
|
+-- venv/                       (Gerado automaticamente)
+-- resultados/                 (Gerado automaticamente)
\end{verbatim}

\subsection{Solução de problemas comuns}

\begin{itemize}
    \item \textbf{Erro "Python não encontrado" ou comando "py" inválido:} O Python não está instalado corretamente. Reinstale marcando a opção "Add to PATH".
    \item \textbf{Erro "Ambiente não encontrado":} A pasta \texttt{venv} pode ter sido deletada. Execute novamente o arquivo \texttt{criar\_ambiente.bat}.
    \item \textbf{Erro "Arquivo não encontrado":} Certifique-se de executar os arquivos \texttt{.bat} sempre a partir da pasta raiz do projeto, e não de dentro das subpastas.
\end{itemize}

\vspace{0.5cm}
\textbf{Nota de portabilidade:} Esta ferramenta é portável. Para usá-la em outro computador, basta copiar a pasta inteira do projeto e rodar o \texttt{criar\_ambiente.bat} no novo local.

\newpage

% \section{Arquitetura de pré-processamento e simulação espacial}
% \label{sec:pre_processamento}
% % Onde se explica como a atratividade, a zona vermelha de penalidade e o Set-Covering são calculados em Python antes do solver.
% A formulação de um problema de localização espacial com interações de mercado (canibalização de demanda) apresenta um desafio inerente: a não-linearidade. Se a atratividade de um novo eletroposto depende de quais outros eletropostos vizinhos serão abertos simultaneamente, a função objetivo original assumiria uma forma bilineal, tornando o problema intratável para solvers exatos em instâncias de grande porte.

Para contornar este obstáculo e garantir uma arquitetura puramente linear (MILP), esta proposta adota uma etapa rigorosa de pré-processamento e simulação espacial. Todo o esforço computacional não-linear é isolado nesta fase, que transforma a topologia urbana em um conjunto de parâmetros fixos que alimentarão o otimizador.

Esta etapa é dividida em quatro fases principais:

\subsection{Geração de candidatos e atratividade gravitacional base}

A área de estudo é inicialmente dividida em uma malha virtual. Para cada quadrante, os Pontos de Interesse (POIs) -- categorizados e ponderados segundo seu potencial de geração de viagens (ex: supermercados, hospitais, shopping centers) -- são mapeados utilizando dados de geolocalização.

O centroide matemático de cada quadrante é calculado e, para garantir o realismo da infraestrutura, aplica-se uma técnica de \textit{Snap to Road} (Ajuste à via). Esta técnica utiliza \textit{Reverse Geocoding} para deslocar o centroide teórico (que poderia recair sobre corpos d'água ou propriedades privadas) para a via pública trafegável mais próxima. Este ponto ajustado torna-se um \textbf{nó candidato} para a instalação de um eletroposto.

Nesta fase, calcula-se a \textbf{Atratividade Base} de cada candidato, que representa a demanda potencial que o local capturaria se operasse como um monopólio isolado.

\subsection{Concorrência existente e decaimento espacial}

Um modelo realista não pode ignorar a infraestrutura de recarga já consolidada na cidade. A presença de um eletroposto concorrente próximo reduz a fração de demanda capturável por um novo candidato.

Para modelar essa interação, implementou-se uma função de \textbf{Decaimento Espacial} baseada na distância real de roteamento (estimada via distância de Haversine com fator de tortuosidade urbana). O modelo estabelece três zonas de impacto para penalizar a atratividade de um candidato:

\begin{itemize}
    \item \textbf{Zona de Impacto Total (Zona Vermelha):} Se um concorrente está demasiadamente próximo (ex: raio inferior a 300 metros), assume-se sobreposição severa de mercado, aplicando-se a penalidade máxima definida.
    \item \textbf{Zona de Transição:} Para distâncias intermediárias (ex: entre 300 e 1500 metros), a penalidade decai linearmente, reduzindo o impacto à medida que os postos se distanciam.
    \item \textbf{Zona Livre:} Se a distância supera o limite máximo de influência (ex: acima de 1500 metros), o concorrente é desconsiderado, assumindo-se impacto nulo.
\end{itemize}

Ao subtrair as penalidades dos concorrentes reais da Atratividade Base (e descontando uma barreira de custo operacional/viabilidade mínima), obtém-se o \textbf{Retorno Líquido Próprio} do candidato.

\subsection{Canibalização potencial e a matriz de cenários}

A inovação central deste pré-processamento reside no tratamento da concorrência entre os próprios nós candidatos. Como ainda não se sabe quais candidatos serão abertos pelo modelo, calcula-se previamente o impacto de \textit{todas as combinações possíveis}.

Para um dado candidato, o algoritmo identifica seus vizinhos potenciais dentro do raio de influência. Para evitar uma explosão combinatória, restringe-se a análise aos $K$ vizinhos mais impactantes (ex: os 4 ou 5 mais próximos). Em seguida, o simulador gera uma tabela exaustiva contendo o retorno líquido do candidato para todos os cenários de operação de seus vizinhos ($2^K$ configurações). 

Este "menu de retornos" pré-calculado é o que permite ao modelo matemático (Seção \ref{sec:modelo}) realizar escolhas de forma puramente linear, apenas selecionando o cenário correspondente à decisão topológica.

\subsection{Mapeamento de cobertura (\textit{Set-Covering})}

A última etapa do pré-processamento garante os insumos para as métricas de qualidade de rede. O modelo exige que uma porcentagem mínima de todos os POIs da região tenha acesso a pelo menos um eletroposto a uma distância caminhável aceitável (ex: 800 metros, ou cerca de 10 minutos de caminhada).

O simulador realiza uma varredura bidirecional:
\begin{enumerate}
    \item \textbf{Cobertura consolidada:} Identifica todos os POIs que já estão a menos da distância caminhável de um eletroposto real existente. Estes são marcados como "atendidos" e isolados, evitando esforços redundantes do modelo.
    \item \textbf{Mapeamento de viabilidade:} Para os POIs restantes (desassistidos), o sistema mapeia quais nós candidatos têm a capacidade geométrica de cobri-los.
\end{enumerate}

Este mapeamento previne o problema estatístico de \textit{dupla contagem} -- garantindo que se dois eletropostos novos cobrirem o mesmo supermercado, a métrica global de cobertura contabilizará o benefício apenas uma vez.

% \section{Modelo matemático de otimização (MILP)}
% \label{sec:modelo}
% % Aqui você irá colar aquele código LaTeX completo com a Função Objetivo, R1, R2, R3 e R4
% \subsection{Conjuntos}

\begin{itemize}
    \item $C$: Conjunto de locais candidatos para instalação de novos eletropostos.
    \item $E$: Conjunto de eletropostos concorrentes já existentes na malha urbana.
    \item $P$: Conjunto total de Pontos de Interesse (POIs) -- hospitais, educação, transporte, comércio.
    \item $P_{\text{cob}} \subseteq P$: Subconjunto de POIs já atendidos (cobertos) por pelo menos um eletroposto existente $e \in E$.
    \item $P_{\text{ncob}} = P \setminus P_{\text{cob}}$: Subconjunto de POIs não atendidos, que dependem das decisões de instalação em $C$.
    \item $V(c)$: Conjunto de vizinhos espaciais do candidato $c$ (dentro do raio crítico de influência).
    \item $\mathcal{S}_c$: Conjunto de todas as combinações possíveis de vizinhos de $c$ que podem ser designados simultaneamente ($\mathcal{S}_c \subseteq 2^{V(c)}$).
\end{itemize}

\subsection{Parâmetros}

\subsubsection{Lucro e Atratividade (pré-computados via simulador espacial)}

\begin{itemize}
    \item $A_c$: Atratividade gravitacional base do candidato $c$ (Score derivado dos POIs próximos).
    \item $B$: Barreira de viabilidade (custo fixo mínimo ou atratividade mínima requerida para operação neutra).
    \item $\Pi_c$: Penalidade aplicada à atratividade do candidato $c$ devido à proximidade (sobreposição de áreas de influência) com eletropostos já existentes $e \in E$.
    
    \item $R_c^{\text{próprio}}$: Retorno líquido do eletroposto $c$ se operasse sem novos vizinhos, mas \textbf{sujeito à concorrência existente}. Definido como:
    \begin{equation}
        R_c^{\text{próprio}} = A_c - B - \Pi_c
    \end{equation}
    
    \item $\Delta R_{c,c'}$: Retorno (demanda) desviado do candidato $c$ para seu vizinho $c' \in V(c)$ quando $c'$ é designado. Representa a canibalização espacial entre as novas opções, calculada com base na distância de decaimento pela malha viária.
    
    \item $R_c(S_c)$: Retorno esperado do candidato $c$ quando exatamente o conjunto de vizinhos $S_c \subseteq V(c)$ está designado. Definido como:
    \begin{equation}
    R_c(S_c) = R_c^{\text{próprio}} - \sum_{c' \in S_c} \Delta R_{c,c'}
    \end{equation}
\end{itemize}

\subsubsection{Qualidade de rede}

\begin{itemize}
    \item $Q_{\min}^{\text{cob}}$: Cobertura mínima da totalidade de POIs requerida na cidade (percentual, ex: 85\%).
    \item $r_{\text{cob}}$: Raio de cobertura máximo aceitável pela malha viária para um POI ser considerado atendido (metros).
\end{itemize}

\subsection{Variáveis de decisão}

\begin{itemize}
    \item $x_c \in \{0,1\}$: Indicadora -- designar a instalação de um eletroposto no candidato $c$.
    \item $y_p \in \{0,1\}$: Indicadora -- POI $p \in P_{\text{ncob}}$ passa a ser coberto por pelo menos um \textit{novo} eletroposto designado.
    \item $z_{c,S_c} \in \{0,1\}$: Indicadora -- a configuração exata de vizinhos \textit{novos} designados para o eletroposto $c$ é o conjunto $S_c$.
\end{itemize}

\subsection{Função objetivo}

\begin{equation}
\label{eq:objetivo_linear}
\max \quad Z = \sum_{c \in C} \sum_{S_c \in \mathcal{S}_c} R_c(S_c) \cdot z_{c,S_c}
\end{equation}

\textit{Nota: O objetivo maximiza o retorno líquido global, selecionando a configuração de concorrência real ($z_{c,S_c}$) que a rede terminará enfrentando após as designações.}

\subsection{Restrições}

\subsubsection{Grupo 1: Linearização e Configuração de Vizinhos}

\hspace{0.5cm} \textbf{R1: Atribuição única de configuração}

\begin{align}
\sum_{S_c \in \mathcal{S}_c} z_{c,S_c} &= x_c \quad \forall c \in C \label{eq:restricao_config}
\end{align}

\hspace{0.5cm} \textbf{R2: Consistência lógica de designações}

Se a configuração $S_c$ for a escolhida ($z_{c,S_c} = 1$), os vizinhos correspondentes devem estar obrigatoriamente designados ou desligados:

\begin{align}
x_{c'} &\geq z_{c,S_c} \quad \forall c \in C, \forall S_c \in \mathcal{S}_c, \forall c' \in S_c \label{eq:restricao_consistencia_in} \\
x_{c'} &\leq 1 - z_{c,S_c} \quad \forall c \in C, \forall S_c \in \mathcal{S}_c, \forall c' \in V(c) \setminus S_c \label{eq:restricao_consistencia_out}
\end{align}

\subsubsection{Grupo 2: Qualidade Global de Rede (Considerando Infraestrutura Existente)}

\hspace{0.5cm} \textbf{R3: Indicador de nova cobertura para POIs desassistidos}

\begin{align}
y_p &\leq \sum_{c \in C: \text{dist}(c,p) \leq r_{\text{cob}}} x_c \quad \forall p \in P_{\text{ncob}} \label{eq:cobertura_poi}
\end{align}

Um POI desassistido $p$ só será contabilizado como coberto ($y_p = 1$) se pelo menos um novo candidato em seu raio de alcance for instalado.

\hspace{0.5cm} \textbf{R4: Meta integrada de cobertura mínima}

\begin{align}
|P_{\text{cob}}| + \sum_{p \in P_{\text{ncob}}} y_p &\geq \lceil Q_{\min}^{\text{cob}} \cdot |P| \rceil \label{eq:meta_cobertura}
\end{align}

A quantidade de POIs cobertos gratuitamente pela infraestrutura existente ($|P_{\text{cob}}|$) somada aos POIs resgatados pela nova infraestrutura deve atender ou superar o volume mínimo exigido pelo planejamento urbano.

% \section{Limitações e extensões futuras}
% \label{sec:extensões}
% % Seção deixada em branco por enquanto, para preencher ao longo do doutorado.
% \input{sections/04_extensões_futuras}

% \end{document}









% \documentclass[a4paper,12pt]{article}

% % --- PACOTES FUNDAMENTAIS ---
% \usepackage[utf8]{inputenc}
% \usepackage[T1]{fontenc}
% \usepackage[brazilian]{babel}
% \usepackage{amsmath}
% \usepackage{amssymb}
% \usepackage{geometry}
% \usepackage{graphicx}
% \usepackage{setspace}
% \usepackage{fancyhdr}
% \usepackage{xcolor} % Para cores personalizadas

% % --- CONFIGURAÇÃO DE LINKS (HYPERREF) ---
% \usepackage{hyperref} 

% % Definição de azul escuro profissional (opcional)
% \definecolor{darkblue}{rgb}{0.0, 0.0, 0.55}

% \hypersetup{
%     colorlinks=true,        % Ativa cores nos links
%     linkcolor=black,        % Cor dos links internos (Sumário, referências). Preto para elegância.
%     filecolor=magenta,      
%     urlcolor=blue,          % Cor de URLs externas (GitHub, LinkedIn)
%     citecolor=black,        % Cor de citações bibliográficas
%     pdftitle={Otimização de Eletropostos (MILP) - Proposta Preliminar},
%     pdfpagemode=FullScreen,
%     linktoc=all             % Torna a linha inteira do sumário clicável
% }

% % --- CONFIGURAÇÃO DE MARGENS ---
% \geometry{left=2.5cm,right=2.5cm,top=2.5cm,bottom=2.5cm}

% \begin{document}

% % --- CAPA ---
% \begin{titlepage}
%     \centering
    
%     % Cabeçalho Institucional
%     {\scshape \Large \textbf{UNIVERSIDADE ESTADUAL DE CAMPINAS (UNICAMP)}}\\[0.5cm]
%     {\large Faculdade de Engenharia Elétrica e de Computação (FEEC)}\\[2.0cm]

%     % Título do Projeto
%     \rule{\linewidth}{0.5mm} \\[0.4cm]
%     { \huge \bfseries Proposta preliminar: \\[0.3cm] Otimização da localização de infraestrutura de recarga para veículos elétricos }\\[0.2cm]
%     { \Large Abordagem via Programação Linear Inteira Mista (MILP) com mitigação de canibalização espacial }\\[0.4cm]
%     \rule{\linewidth}{0.5mm} \\[2.5cm]
    
%     % Informações dos Pesquisadores
%     \begin{minipage}{0.8\textwidth}
%         \centering
%         \textbf{Pesquisadores:}\\[0.4cm]
%         {\Large Núbia Sidene Almeida das Virgens}\\[0.2em]
%         \textit{Revisão e co-desenvolvimento da etapa de otimização}\\[0.8cm]

%         {\Large Carlos Julian Muñoz Quiroga}\\[0.2em]
%         \textit{Doutorando em Engenharia Elétrica e de Computação}\\
%     \end{minipage}
    
%     \vspace{2cm}
    
%     % Links (Opcional - Pode remover se o repositório for privado)
%     \textbf{Repositório do protótipo e dashboard:}\\
%     \href{https://github.com/CAMUZQUI/designacao-eletropostos-milp}{designacao-eletropostos-milp}
    
%     \vfill
    
%     % Rodapé
%     {\large Campinas, SP}\\
%     {\large \today}
    
% \end{titlepage}

% % --- SUMÁRIO ---
% \newpage
% \thispagestyle{empty} 

% \tableofcontents 

% \newpage

% % --- CONTEÚDO MODULAR ---

% \section{Introdução e contexto do problema}
% \label{sec:introducao}
% Este guia orienta a configuração e utilização da ferramenta de otimização de forma direta e autônoma.

\subsection{Pré-requisitos do sistema}

Para executar esta ferramenta, seu ambiente deve atender a dois requisitos fundamentais:

\subsubsection{Python 3.10}
Verifique se o Python está instalado abrindo o \textit{Prompt de Comando} (cmd) e digitando:

\begin{verbatim}
py --list
\end{verbatim}

\textbf{Como interpretar:}
\begin{itemize}
    \item Se aparecer \texttt{ - 3.10.X} (ou superior), o Python está pronto.
    \item Se der erro, baixe o instalador em \texttt{python.org} e marque a opção "Add to PATH" durante a instalação.
\end{itemize}

\subsubsection{IBM ILOG CPLEX Optimization Studio}
O modelo utiliza a biblioteca \texttt{docplex}, que requer o motor de otimização da IBM instalado no sistema.

\begin{itemize}
    \item É obrigatório ter uma instalação do CPLEX com licença acadêmica ou comercial.
    \item Não utilize a versão "Community Edition", ela possui limites restritos de variáveis e restrições que são insuficientes para a complexidade deste case, impedindo a resolução.
\end{itemize}


\subsection{Configuração do ambiente}

Este passo é necessário apenas na primeira vez que você for usar a ferramenta. Localize e execute o arquivo \texttt{criar\_ambiente.bat} na pasta raiz do projeto.
\begin{itemize}
    \item Uma janela do terminal se abrirá e instalará automaticamente as dependências.
    \item Aguarde até ver a mensagem "AMBIENTE CRIADO COM SUCESSO" e feche a janela.
    \item Se a janela fechar instantaneamente sem mensagens de sucesso, verifique sua instalação do Python.
\end{itemize}

\subsection{Execução da ferramenta}

Para utilizar o otimizador no dia a dia:

\begin{enumerate}
    \item Execute o arquivo \texttt{executar\_solver.bat} na pasta raiz.
    \item Uma janela de seleção de arquivos abrirá automaticamente. Selecione a planilha Excel com os dados de entrada.
    \item Acompanhe o progresso no terminal. O solver processará os dados e buscará a melhor solução.
    \item Ao finalizar, a ferramenta abrirá automaticamente a planilha de resultados gerada na pasta \texttt{resultados}.
\end{enumerate}

\vspace{0.3cm}

\textbf{Controle de Execução (Importante):}

O solver foi configurado para buscar a solução ótima sem um limite de tempo pré-fixado.
\begin{itemize}
    \item \textbf{Interrupção manual:} Você pode interromper o processo a qualquer momento pressionando o atalho \texttt{Ctrl + C} na janela do terminal.
    \item \textbf{Salvamento automático:} Ao interromper o processo, o programa salvará automaticamente a melhor solução encontrada até aquele instante, garantindo que nenhum progresso válido seja perdido.
\end{itemize}


\subsection{Estrutura de arquivos}

A organização das pastas do projeto segue o padrão abaixo. Mantenha essa estrutura para garantir o funcionamento dos scripts automáticos.

\begin{verbatim}
desafio-otimizacao-suzano/
|
+-- criar_ambiente.bat          (Configuracao inicial)
+-- executar_solver.bat         (Execucao diaria)
+-- requirements.txt            (Lista de dependencias)
|
+-- code/                       (Codigos-fonte Python)
|   +-- modelo.py
|   +-- gestor_dados.py
|   +-- exportar_resultados.py
|
+-- venv/                       (Gerado automaticamente)
+-- resultados/                 (Gerado automaticamente)
\end{verbatim}

\subsection{Solução de problemas comuns}

\begin{itemize}
    \item \textbf{Erro "Python não encontrado" ou comando "py" inválido:} O Python não está instalado corretamente. Reinstale marcando a opção "Add to PATH".
    \item \textbf{Erro "Ambiente não encontrado":} A pasta \texttt{venv} pode ter sido deletada. Execute novamente o arquivo \texttt{criar\_ambiente.bat}.
    \item \textbf{Erro "Arquivo não encontrado":} Certifique-se de executar os arquivos \texttt{.bat} sempre a partir da pasta raiz do projeto, e não de dentro das subpastas.
\end{itemize}

\vspace{0.5cm}
\textbf{Nota de portabilidade:} Esta ferramenta é portável. Para usá-la em outro computador, basta copiar a pasta inteira do projeto e rodar o \texttt{criar\_ambiente.bat} no novo local.

\newpage

% \newpage
% \section{Arquitetura de pré-processamento e simulação espacial}
% \label{sec:pre_processamento}
% % Onde se explica como a atratividade, a zona vermelha de penalidade e o Set-Covering são calculados em Python antes do solver.
% A formulação de um problema de localização espacial com interações de mercado (canibalização de demanda) apresenta um desafio inerente: a não-linearidade. Se a atratividade de um novo eletroposto depende de quais outros eletropostos vizinhos serão abertos simultaneamente, a função objetivo original assumiria uma forma bilineal, tornando o problema intratável para solvers exatos em instâncias de grande porte.

Para contornar este obstáculo e garantir uma arquitetura puramente linear (MILP), esta proposta adota uma etapa rigorosa de pré-processamento e simulação espacial. Todo o esforço computacional não-linear é isolado nesta fase, que transforma a topologia urbana em um conjunto de parâmetros fixos que alimentarão o otimizador.

Esta etapa é dividida em quatro fases principais:

\subsection{Geração de candidatos e atratividade gravitacional base}

A área de estudo é inicialmente dividida em uma malha virtual. Para cada quadrante, os Pontos de Interesse (POIs) -- categorizados e ponderados segundo seu potencial de geração de viagens (ex: supermercados, hospitais, shopping centers) -- são mapeados utilizando dados de geolocalização.

O centroide matemático de cada quadrante é calculado e, para garantir o realismo da infraestrutura, aplica-se uma técnica de \textit{Snap to Road} (Ajuste à via). Esta técnica utiliza \textit{Reverse Geocoding} para deslocar o centroide teórico (que poderia recair sobre corpos d'água ou propriedades privadas) para a via pública trafegável mais próxima. Este ponto ajustado torna-se um \textbf{nó candidato} para a instalação de um eletroposto.

Nesta fase, calcula-se a \textbf{Atratividade Base} de cada candidato, que representa a demanda potencial que o local capturaria se operasse como um monopólio isolado.

\subsection{Concorrência existente e decaimento espacial}

Um modelo realista não pode ignorar a infraestrutura de recarga já consolidada na cidade. A presença de um eletroposto concorrente próximo reduz a fração de demanda capturável por um novo candidato.

Para modelar essa interação, implementou-se uma função de \textbf{Decaimento Espacial} baseada na distância real de roteamento (estimada via distância de Haversine com fator de tortuosidade urbana). O modelo estabelece três zonas de impacto para penalizar a atratividade de um candidato:

\begin{itemize}
    \item \textbf{Zona de Impacto Total (Zona Vermelha):} Se um concorrente está demasiadamente próximo (ex: raio inferior a 300 metros), assume-se sobreposição severa de mercado, aplicando-se a penalidade máxima definida.
    \item \textbf{Zona de Transição:} Para distâncias intermediárias (ex: entre 300 e 1500 metros), a penalidade decai linearmente, reduzindo o impacto à medida que os postos se distanciam.
    \item \textbf{Zona Livre:} Se a distância supera o limite máximo de influência (ex: acima de 1500 metros), o concorrente é desconsiderado, assumindo-se impacto nulo.
\end{itemize}

Ao subtrair as penalidades dos concorrentes reais da Atratividade Base (e descontando uma barreira de custo operacional/viabilidade mínima), obtém-se o \textbf{Retorno Líquido Próprio} do candidato.

\subsection{Canibalização potencial e a matriz de cenários}

A inovação central deste pré-processamento reside no tratamento da concorrência entre os próprios nós candidatos. Como ainda não se sabe quais candidatos serão abertos pelo modelo, calcula-se previamente o impacto de \textit{todas as combinações possíveis}.

Para um dado candidato, o algoritmo identifica seus vizinhos potenciais dentro do raio de influência. Para evitar uma explosão combinatória, restringe-se a análise aos $K$ vizinhos mais impactantes (ex: os 4 ou 5 mais próximos). Em seguida, o simulador gera uma tabela exaustiva contendo o retorno líquido do candidato para todos os cenários de operação de seus vizinhos ($2^K$ configurações). 

Este "menu de retornos" pré-calculado é o que permite ao modelo matemático (Seção \ref{sec:modelo}) realizar escolhas de forma puramente linear, apenas selecionando o cenário correspondente à decisão topológica.

\subsection{Mapeamento de cobertura (\textit{Set-Covering})}

A última etapa do pré-processamento garante os insumos para as métricas de qualidade de rede. O modelo exige que uma porcentagem mínima de todos os POIs da região tenha acesso a pelo menos um eletroposto a uma distância caminhável aceitável (ex: 800 metros, ou cerca de 10 minutos de caminhada).

O simulador realiza uma varredura bidirecional:
\begin{enumerate}
    \item \textbf{Cobertura consolidada:} Identifica todos os POIs que já estão a menos da distância caminhável de um eletroposto real existente. Estes são marcados como "atendidos" e isolados, evitando esforços redundantes do modelo.
    \item \textbf{Mapeamento de viabilidade:} Para os POIs restantes (desassistidos), o sistema mapeia quais nós candidatos têm a capacidade geométrica de cobri-los.
\end{enumerate}

Este mapeamento previne o problema estatístico de \textit{dupla contagem} -- garantindo que se dois eletropostos novos cobrirem o mesmo supermercado, a métrica global de cobertura contabilizará o benefício apenas uma vez.

% \newpage
% \section{Modelo matemático de otimização (MILP)}
% \label{sec:modelo}
% % Aqui você irá colar aquele código LaTeX completo com a Função Objetivo, R1, R2, R3 e R4
% \subsection{Conjuntos}

\begin{itemize}
    \item $C$: Conjunto de locais candidatos para instalação de novos eletropostos.
    \item $E$: Conjunto de eletropostos concorrentes já existentes na malha urbana.
    \item $P$: Conjunto total de Pontos de Interesse (POIs) -- hospitais, educação, transporte, comércio.
    \item $P_{\text{cob}} \subseteq P$: Subconjunto de POIs já atendidos (cobertos) por pelo menos um eletroposto existente $e \in E$.
    \item $P_{\text{ncob}} = P \setminus P_{\text{cob}}$: Subconjunto de POIs não atendidos, que dependem das decisões de instalação em $C$.
    \item $V(c)$: Conjunto de vizinhos espaciais do candidato $c$ (dentro do raio crítico de influência).
    \item $\mathcal{S}_c$: Conjunto de todas as combinações possíveis de vizinhos de $c$ que podem ser designados simultaneamente ($\mathcal{S}_c \subseteq 2^{V(c)}$).
\end{itemize}

\subsection{Parâmetros}

\subsubsection{Lucro e Atratividade (pré-computados via simulador espacial)}

\begin{itemize}
    \item $A_c$: Atratividade gravitacional base do candidato $c$ (Score derivado dos POIs próximos).
    \item $B$: Barreira de viabilidade (custo fixo mínimo ou atratividade mínima requerida para operação neutra).
    \item $\Pi_c$: Penalidade aplicada à atratividade do candidato $c$ devido à proximidade (sobreposição de áreas de influência) com eletropostos já existentes $e \in E$.
    
    \item $R_c^{\text{próprio}}$: Retorno líquido do eletroposto $c$ se operasse sem novos vizinhos, mas \textbf{sujeito à concorrência existente}. Definido como:
    \begin{equation}
        R_c^{\text{próprio}} = A_c - B - \Pi_c
    \end{equation}
    
    \item $\Delta R_{c,c'}$: Retorno (demanda) desviado do candidato $c$ para seu vizinho $c' \in V(c)$ quando $c'$ é designado. Representa a canibalização espacial entre as novas opções, calculada com base na distância de decaimento pela malha viária.
    
    \item $R_c(S_c)$: Retorno esperado do candidato $c$ quando exatamente o conjunto de vizinhos $S_c \subseteq V(c)$ está designado. Definido como:
    \begin{equation}
    R_c(S_c) = R_c^{\text{próprio}} - \sum_{c' \in S_c} \Delta R_{c,c'}
    \end{equation}
\end{itemize}

\subsubsection{Qualidade de rede}

\begin{itemize}
    \item $Q_{\min}^{\text{cob}}$: Cobertura mínima da totalidade de POIs requerida na cidade (percentual, ex: 85\%).
    \item $r_{\text{cob}}$: Raio de cobertura máximo aceitável pela malha viária para um POI ser considerado atendido (metros).
\end{itemize}

\subsection{Variáveis de decisão}

\begin{itemize}
    \item $x_c \in \{0,1\}$: Indicadora -- designar a instalação de um eletroposto no candidato $c$.
    \item $y_p \in \{0,1\}$: Indicadora -- POI $p \in P_{\text{ncob}}$ passa a ser coberto por pelo menos um \textit{novo} eletroposto designado.
    \item $z_{c,S_c} \in \{0,1\}$: Indicadora -- a configuração exata de vizinhos \textit{novos} designados para o eletroposto $c$ é o conjunto $S_c$.
\end{itemize}

\subsection{Função objetivo}

\begin{equation}
\label{eq:objetivo_linear}
\max \quad Z = \sum_{c \in C} \sum_{S_c \in \mathcal{S}_c} R_c(S_c) \cdot z_{c,S_c}
\end{equation}

\textit{Nota: O objetivo maximiza o retorno líquido global, selecionando a configuração de concorrência real ($z_{c,S_c}$) que a rede terminará enfrentando após as designações.}

\subsection{Restrições}

\subsubsection{Grupo 1: Linearização e Configuração de Vizinhos}

\hspace{0.5cm} \textbf{R1: Atribuição única de configuração}

\begin{align}
\sum_{S_c \in \mathcal{S}_c} z_{c,S_c} &= x_c \quad \forall c \in C \label{eq:restricao_config}
\end{align}

\hspace{0.5cm} \textbf{R2: Consistência lógica de designações}

Se a configuração $S_c$ for a escolhida ($z_{c,S_c} = 1$), os vizinhos correspondentes devem estar obrigatoriamente designados ou desligados:

\begin{align}
x_{c'} &\geq z_{c,S_c} \quad \forall c \in C, \forall S_c \in \mathcal{S}_c, \forall c' \in S_c \label{eq:restricao_consistencia_in} \\
x_{c'} &\leq 1 - z_{c,S_c} \quad \forall c \in C, \forall S_c \in \mathcal{S}_c, \forall c' \in V(c) \setminus S_c \label{eq:restricao_consistencia_out}
\end{align}

\subsubsection{Grupo 2: Qualidade Global de Rede (Considerando Infraestrutura Existente)}

\hspace{0.5cm} \textbf{R3: Indicador de nova cobertura para POIs desassistidos}

\begin{align}
y_p &\leq \sum_{c \in C: \text{dist}(c,p) \leq r_{\text{cob}}} x_c \quad \forall p \in P_{\text{ncob}} \label{eq:cobertura_poi}
\end{align}

Um POI desassistido $p$ só será contabilizado como coberto ($y_p = 1$) se pelo menos um novo candidato em seu raio de alcance for instalado.

\hspace{0.5cm} \textbf{R4: Meta integrada de cobertura mínima}

\begin{align}
|P_{\text{cob}}| + \sum_{p \in P_{\text{ncob}}} y_p &\geq \lceil Q_{\min}^{\text{cob}} \cdot |P| \rceil \label{eq:meta_cobertura}
\end{align}

A quantidade de POIs cobertos gratuitamente pela infraestrutura existente ($|P_{\text{cob}}|$) somada aos POIs resgatados pela nova infraestrutura deve atender ou superar o volume mínimo exigido pelo planejamento urbano.

% \newpage
% \section{Limitações e extensões futuras}
% \label{sec:extensões}
% % Seção deixada em branco por enquanto, para preencher ao longo do doutorado.
% \input{sections/04_extensões_futuras}

% \end{document}









\documentclass[a4paper,12pt]{article}

% --- PACOTES FUNDAMENTAIS ---
\usepackage[utf8]{inputenc}
\usepackage[T1]{fontenc}
\usepackage[brazilian]{babel}
\usepackage{amsmath}
\usepackage{amssymb}
\usepackage{geometry}
\usepackage{graphicx}
\usepackage{setspace}
\usepackage{fancyhdr}
\usepackage{xcolor} % Para cores personalizadas

% --- CONFIGURAÇÃO DE LINKS (HYPERREF) ---
\usepackage{hyperref} 

% Definição de azul escuro profissional (opcional)
\definecolor{darkblue}{rgb}{0.0, 0.0, 0.55}

\hypersetup{
    colorlinks=true,        % Ativa cores nos links
    linkcolor=black,        % Cor dos links internos (Sumário, referências). Preto para elegância.
    filecolor=magenta,      
    urlcolor=blue,          % Cor de URLs externas (GitHub, LinkedIn)
    citecolor=black,        % Cor de citações bibliográficas
    pdftitle={Otimização de Eletropostos (MILP) - Proposta Preliminar},
    pdfpagemode=FullScreen,
    linktoc=all             % Torna a linha inteira do sumário clicável
}

% --- CONFIGURAÇÃO DE MARGENS ---
\geometry{left=2.5cm,right=2.5cm,top=2.5cm,bottom=2.5cm}

\begin{document}

% --- CAPA ---
\begin{titlepage}
    \centering
    
    % Cabeçalho Institucional
    {\scshape \Large \textbf{UNIVERSIDADE ESTADUAL DE CAMPINAS (UNICAMP)}}\\[0.5cm]
    {\large Faculdade de Engenharia Elétrica e de Computação (FEEC)}\\[2.0cm]

    % Título do Projeto
    \rule{\linewidth}{0.5mm} \\[0.4cm]
    { \huge \bfseries Proposta preliminar: \\[0.3cm] Otimização da localização de infraestrutura de recarga para veículos elétricos }\\[0.2cm]
    { \Large Abordagem via Programação Linear Inteira Mista (MILP) com mitigação de canibalização espacial }\\[0.4cm]
    \rule{\linewidth}{0.5mm} \\[2.0cm]
    
    % Informações dos Pesquisadores e Orientador
    \begin{minipage}{0.8\textwidth}
        \centering
        \textbf{Pesquisadores:}\\[0.4cm]
        {\Large Núbia Sidene Almeida das Virgens}\\[0.2em]
        \textit{Doutoranda em Engenharia Elétrica e de Computação}\\[0.6cm]

        {\Large Carlos Julian Muñoz Quiroga}\\[0.2em]
        \textit{Doutorando em Engenharia Elétrica e de Computação}\\[0.8cm]

        \textbf{Orientado por:}\\[0.2em]
        {\Large Prof. Dr. Romis Ribeiro de Faissol Attux}\\
    \end{minipage}
    
    \vspace{1.5cm}
    
    % Links (Opcional - Pode remover se o repositório for privado)
    \textbf{Repositório do protótipo e modelo:}\\
    \href{https://github.com/Carmuzqui/Primeira-proposta-MILP}{Primeira-proposta-MILP}
    
    \vfill
    
    % Rodapé
    {\large Campinas, SP}\\
    {\large \today}
    
\end{titlepage}

% --- SUMÁRIO ---
\newpage
\thispagestyle{empty} 

\tableofcontents 

\newpage

% --- CONTEÚDO MODULAR ---

\section{Introdução e contexto do problema}
\label{sec:introducao}
Este guia orienta a configuração e utilização da ferramenta de otimização de forma direta e autônoma.

\subsection{Pré-requisitos do sistema}

Para executar esta ferramenta, seu ambiente deve atender a dois requisitos fundamentais:

\subsubsection{Python 3.10}
Verifique se o Python está instalado abrindo o \textit{Prompt de Comando} (cmd) e digitando:

\begin{verbatim}
py --list
\end{verbatim}

\textbf{Como interpretar:}
\begin{itemize}
    \item Se aparecer \texttt{ - 3.10.X} (ou superior), o Python está pronto.
    \item Se der erro, baixe o instalador em \texttt{python.org} e marque a opção "Add to PATH" durante a instalação.
\end{itemize}

\subsubsection{IBM ILOG CPLEX Optimization Studio}
O modelo utiliza a biblioteca \texttt{docplex}, que requer o motor de otimização da IBM instalado no sistema.

\begin{itemize}
    \item É obrigatório ter uma instalação do CPLEX com licença acadêmica ou comercial.
    \item Não utilize a versão "Community Edition", ela possui limites restritos de variáveis e restrições que são insuficientes para a complexidade deste case, impedindo a resolução.
\end{itemize}


\subsection{Configuração do ambiente}

Este passo é necessário apenas na primeira vez que você for usar a ferramenta. Localize e execute o arquivo \texttt{criar\_ambiente.bat} na pasta raiz do projeto.
\begin{itemize}
    \item Uma janela do terminal se abrirá e instalará automaticamente as dependências.
    \item Aguarde até ver a mensagem "AMBIENTE CRIADO COM SUCESSO" e feche a janela.
    \item Se a janela fechar instantaneamente sem mensagens de sucesso, verifique sua instalação do Python.
\end{itemize}

\subsection{Execução da ferramenta}

Para utilizar o otimizador no dia a dia:

\begin{enumerate}
    \item Execute o arquivo \texttt{executar\_solver.bat} na pasta raiz.
    \item Uma janela de seleção de arquivos abrirá automaticamente. Selecione a planilha Excel com os dados de entrada.
    \item Acompanhe o progresso no terminal. O solver processará os dados e buscará a melhor solução.
    \item Ao finalizar, a ferramenta abrirá automaticamente a planilha de resultados gerada na pasta \texttt{resultados}.
\end{enumerate}

\vspace{0.3cm}

\textbf{Controle de Execução (Importante):}

O solver foi configurado para buscar a solução ótima sem um limite de tempo pré-fixado.
\begin{itemize}
    \item \textbf{Interrupção manual:} Você pode interromper o processo a qualquer momento pressionando o atalho \texttt{Ctrl + C} na janela do terminal.
    \item \textbf{Salvamento automático:} Ao interromper o processo, o programa salvará automaticamente a melhor solução encontrada até aquele instante, garantindo que nenhum progresso válido seja perdido.
\end{itemize}


\subsection{Estrutura de arquivos}

A organização das pastas do projeto segue o padrão abaixo. Mantenha essa estrutura para garantir o funcionamento dos scripts automáticos.

\begin{verbatim}
desafio-otimizacao-suzano/
|
+-- criar_ambiente.bat          (Configuracao inicial)
+-- executar_solver.bat         (Execucao diaria)
+-- requirements.txt            (Lista de dependencias)
|
+-- code/                       (Codigos-fonte Python)
|   +-- modelo.py
|   +-- gestor_dados.py
|   +-- exportar_resultados.py
|
+-- venv/                       (Gerado automaticamente)
+-- resultados/                 (Gerado automaticamente)
\end{verbatim}

\subsection{Solução de problemas comuns}

\begin{itemize}
    \item \textbf{Erro "Python não encontrado" ou comando "py" inválido:} O Python não está instalado corretamente. Reinstale marcando a opção "Add to PATH".
    \item \textbf{Erro "Ambiente não encontrado":} A pasta \texttt{venv} pode ter sido deletada. Execute novamente o arquivo \texttt{criar\_ambiente.bat}.
    \item \textbf{Erro "Arquivo não encontrado":} Certifique-se de executar os arquivos \texttt{.bat} sempre a partir da pasta raiz do projeto, e não de dentro das subpastas.
\end{itemize}

\vspace{0.5cm}
\textbf{Nota de portabilidade:} Esta ferramenta é portável. Para usá-la em outro computador, basta copiar a pasta inteira do projeto e rodar o \texttt{criar\_ambiente.bat} no novo local.

\newpage

\newpage
\section{Arquitetura de pré-processamento e simulação espacial}
\label{sec:pre_processamento}
% Onde se explica como a atratividade, a zona vermelha de penalidade e o Set-Covering são calculados em Python antes do solver.
A formulação de um problema de localização espacial com interações de mercado (canibalização de demanda) apresenta um desafio inerente: a não-linearidade. Se a atratividade de um novo eletroposto depende de quais outros eletropostos vizinhos serão abertos simultaneamente, a função objetivo original assumiria uma forma bilineal, tornando o problema intratável para solvers exatos em instâncias de grande porte.

Para contornar este obstáculo e garantir uma arquitetura puramente linear (MILP), esta proposta adota uma etapa rigorosa de pré-processamento e simulação espacial. Todo o esforço computacional não-linear é isolado nesta fase, que transforma a topologia urbana em um conjunto de parâmetros fixos que alimentarão o otimizador.

Esta etapa é dividida em quatro fases principais:

\subsection{Geração de candidatos e atratividade gravitacional base}

A área de estudo é inicialmente dividida em uma malha virtual. Para cada quadrante, os Pontos de Interesse (POIs) -- categorizados e ponderados segundo seu potencial de geração de viagens (ex: supermercados, hospitais, shopping centers) -- são mapeados utilizando dados de geolocalização.

O centroide matemático de cada quadrante é calculado e, para garantir o realismo da infraestrutura, aplica-se uma técnica de \textit{Snap to Road} (Ajuste à via). Esta técnica utiliza \textit{Reverse Geocoding} para deslocar o centroide teórico (que poderia recair sobre corpos d'água ou propriedades privadas) para a via pública trafegável mais próxima. Este ponto ajustado torna-se um \textbf{nó candidato} para a instalação de um eletroposto.

Nesta fase, calcula-se a \textbf{Atratividade Base} de cada candidato, que representa a demanda potencial que o local capturaria se operasse como um monopólio isolado.

\subsection{Concorrência existente e decaimento espacial}

Um modelo realista não pode ignorar a infraestrutura de recarga já consolidada na cidade. A presença de um eletroposto concorrente próximo reduz a fração de demanda capturável por um novo candidato.

Para modelar essa interação, implementou-se uma função de \textbf{Decaimento Espacial} baseada na distância real de roteamento (estimada via distância de Haversine com fator de tortuosidade urbana). O modelo estabelece três zonas de impacto para penalizar a atratividade de um candidato:

\begin{itemize}
    \item \textbf{Zona de Impacto Total (Zona Vermelha):} Se um concorrente está demasiadamente próximo (ex: raio inferior a 300 metros), assume-se sobreposição severa de mercado, aplicando-se a penalidade máxima definida.
    \item \textbf{Zona de Transição:} Para distâncias intermediárias (ex: entre 300 e 1500 metros), a penalidade decai linearmente, reduzindo o impacto à medida que os postos se distanciam.
    \item \textbf{Zona Livre:} Se a distância supera o limite máximo de influência (ex: acima de 1500 metros), o concorrente é desconsiderado, assumindo-se impacto nulo.
\end{itemize}

Ao subtrair as penalidades dos concorrentes reais da Atratividade Base (e descontando uma barreira de custo operacional/viabilidade mínima), obtém-se o \textbf{Retorno Líquido Próprio} do candidato.

\subsection{Canibalização potencial e a matriz de cenários}

A inovação central deste pré-processamento reside no tratamento da concorrência entre os próprios nós candidatos. Como ainda não se sabe quais candidatos serão abertos pelo modelo, calcula-se previamente o impacto de \textit{todas as combinações possíveis}.

Para um dado candidato, o algoritmo identifica seus vizinhos potenciais dentro do raio de influência. Para evitar uma explosão combinatória, restringe-se a análise aos $K$ vizinhos mais impactantes (ex: os 4 ou 5 mais próximos). Em seguida, o simulador gera uma tabela exaustiva contendo o retorno líquido do candidato para todos os cenários de operação de seus vizinhos ($2^K$ configurações). 

Este "menu de retornos" pré-calculado é o que permite ao modelo matemático (Seção \ref{sec:modelo}) realizar escolhas de forma puramente linear, apenas selecionando o cenário correspondente à decisão topológica.

\subsection{Mapeamento de cobertura (\textit{Set-Covering})}

A última etapa do pré-processamento garante os insumos para as métricas de qualidade de rede. O modelo exige que uma porcentagem mínima de todos os POIs da região tenha acesso a pelo menos um eletroposto a uma distância caminhável aceitável (ex: 800 metros, ou cerca de 10 minutos de caminhada).

O simulador realiza uma varredura bidirecional:
\begin{enumerate}
    \item \textbf{Cobertura consolidada:} Identifica todos os POIs que já estão a menos da distância caminhável de um eletroposto real existente. Estes são marcados como "atendidos" e isolados, evitando esforços redundantes do modelo.
    \item \textbf{Mapeamento de viabilidade:} Para os POIs restantes (desassistidos), o sistema mapeia quais nós candidatos têm a capacidade geométrica de cobri-los.
\end{enumerate}

Este mapeamento previne o problema estatístico de \textit{dupla contagem} -- garantindo que se dois eletropostos novos cobrirem o mesmo supermercado, a métrica global de cobertura contabilizará o benefício apenas uma vez.

\newpage
\section{Modelo matemático de otimização (MILP)}
\label{sec:modelo}
% Aqui você irá colar aquele código LaTeX completo com a Função Objetivo, R1, R2, R3 e R4
\subsection{Conjuntos}

\begin{itemize}
    \item $C$: Conjunto de locais candidatos para instalação de novos eletropostos.
    \item $E$: Conjunto de eletropostos concorrentes já existentes na malha urbana.
    \item $P$: Conjunto total de Pontos de Interesse (POIs) -- hospitais, educação, transporte, comércio.
    \item $P_{\text{cob}} \subseteq P$: Subconjunto de POIs já atendidos (cobertos) por pelo menos um eletroposto existente $e \in E$.
    \item $P_{\text{ncob}} = P \setminus P_{\text{cob}}$: Subconjunto de POIs não atendidos, que dependem das decisões de instalação em $C$.
    \item $V(c)$: Conjunto de vizinhos espaciais do candidato $c$ (dentro do raio crítico de influência).
    \item $\mathcal{S}_c$: Conjunto de todas as combinações possíveis de vizinhos de $c$ que podem ser designados simultaneamente ($\mathcal{S}_c \subseteq 2^{V(c)}$).
\end{itemize}

\subsection{Parâmetros}

\subsubsection{Lucro e Atratividade (pré-computados via simulador espacial)}

\begin{itemize}
    \item $A_c$: Atratividade gravitacional base do candidato $c$ (Score derivado dos POIs próximos).
    \item $B$: Barreira de viabilidade (custo fixo mínimo ou atratividade mínima requerida para operação neutra).
    \item $\Pi_c$: Penalidade aplicada à atratividade do candidato $c$ devido à proximidade (sobreposição de áreas de influência) com eletropostos já existentes $e \in E$.
    
    \item $R_c^{\text{próprio}}$: Retorno líquido do eletroposto $c$ se operasse sem novos vizinhos, mas \textbf{sujeito à concorrência existente}. Definido como:
    \begin{equation}
        R_c^{\text{próprio}} = A_c - B - \Pi_c
    \end{equation}
    
    \item $\Delta R_{c,c'}$: Retorno (demanda) desviado do candidato $c$ para seu vizinho $c' \in V(c)$ quando $c'$ é designado. Representa a canibalização espacial entre as novas opções, calculada com base na distância de decaimento pela malha viária.
    
    \item $R_c(S_c)$: Retorno esperado do candidato $c$ quando exatamente o conjunto de vizinhos $S_c \subseteq V(c)$ está designado. Definido como:
    \begin{equation}
    R_c(S_c) = R_c^{\text{próprio}} - \sum_{c' \in S_c} \Delta R_{c,c'}
    \end{equation}
\end{itemize}

\subsubsection{Qualidade de rede}

\begin{itemize}
    \item $Q_{\min}^{\text{cob}}$: Cobertura mínima da totalidade de POIs requerida na cidade (percentual, ex: 85\%).
    \item $r_{\text{cob}}$: Raio de cobertura máximo aceitável pela malha viária para um POI ser considerado atendido (metros).
\end{itemize}

\subsection{Variáveis de decisão}

\begin{itemize}
    \item $x_c \in \{0,1\}$: Indicadora -- designar a instalação de um eletroposto no candidato $c$.
    \item $y_p \in \{0,1\}$: Indicadora -- POI $p \in P_{\text{ncob}}$ passa a ser coberto por pelo menos um \textit{novo} eletroposto designado.
    \item $z_{c,S_c} \in \{0,1\}$: Indicadora -- a configuração exata de vizinhos \textit{novos} designados para o eletroposto $c$ é o conjunto $S_c$.
\end{itemize}

\subsection{Função objetivo}

\begin{equation}
\label{eq:objetivo_linear}
\max \quad Z = \sum_{c \in C} \sum_{S_c \in \mathcal{S}_c} R_c(S_c) \cdot z_{c,S_c}
\end{equation}

\textit{Nota: O objetivo maximiza o retorno líquido global, selecionando a configuração de concorrência real ($z_{c,S_c}$) que a rede terminará enfrentando após as designações.}

\subsection{Restrições}

\subsubsection{Grupo 1: Linearização e Configuração de Vizinhos}

\hspace{0.5cm} \textbf{R1: Atribuição única de configuração}

\begin{align}
\sum_{S_c \in \mathcal{S}_c} z_{c,S_c} &= x_c \quad \forall c \in C \label{eq:restricao_config}
\end{align}

\hspace{0.5cm} \textbf{R2: Consistência lógica de designações}

Se a configuração $S_c$ for a escolhida ($z_{c,S_c} = 1$), os vizinhos correspondentes devem estar obrigatoriamente designados ou desligados:

\begin{align}
x_{c'} &\geq z_{c,S_c} \quad \forall c \in C, \forall S_c \in \mathcal{S}_c, \forall c' \in S_c \label{eq:restricao_consistencia_in} \\
x_{c'} &\leq 1 - z_{c,S_c} \quad \forall c \in C, \forall S_c \in \mathcal{S}_c, \forall c' \in V(c) \setminus S_c \label{eq:restricao_consistencia_out}
\end{align}

\subsubsection{Grupo 2: Qualidade Global de Rede (Considerando Infraestrutura Existente)}

\hspace{0.5cm} \textbf{R3: Indicador de nova cobertura para POIs desassistidos}

\begin{align}
y_p &\leq \sum_{c \in C: \text{dist}(c,p) \leq r_{\text{cob}}} x_c \quad \forall p \in P_{\text{ncob}} \label{eq:cobertura_poi}
\end{align}

Um POI desassistido $p$ só será contabilizado como coberto ($y_p = 1$) se pelo menos um novo candidato em seu raio de alcance for instalado.

\hspace{0.5cm} \textbf{R4: Meta integrada de cobertura mínima}

\begin{align}
|P_{\text{cob}}| + \sum_{p \in P_{\text{ncob}}} y_p &\geq \lceil Q_{\min}^{\text{cob}} \cdot |P| \rceil \label{eq:meta_cobertura}
\end{align}

A quantidade de POIs cobertos gratuitamente pela infraestrutura existente ($|P_{\text{cob}}|$) somada aos POIs resgatados pela nova infraestrutura deve atender ou superar o volume mínimo exigido pelo planejamento urbano.

\newpage
\section{Limitações e extensões futuras}
\label{sec:extensoes}
% Seção deixada em branco por enquanto, para preencher ao longo do doutorado.
A arquitetura proposta neste relatório estabelece uma base matemática sólida e computacionalmente eficiente para a localização de infraestrutura de recarga de veículos elétricos (EVCS). O isolamento das não-linearidades em uma etapa prévia de simulação espacial permite que o otimizador MILP opere com máxima performance. 

Contudo, como toda abstração da realidade urbana, o modelo atual apresenta premissas simplificadoras. Para as próximas etapas desta pesquisa de doutorado, propõe-se a expansão do modelo através das seguintes frentes de desenvolvimento:

\subsection{Refinamento dos parâmetros de concorrência e demanda}

Atualmente, o decaimento espacial e a canibalização são governados estritamente pela distância viária. Embora robusta, essa abordagem pode ser enriquecida:
\begin{itemize}
    \item \textbf{Heterogeneidade da infraestrutura:} Incorporar metadados detalhados dos eletropostos (ex: quantidade de conectores, potência máxima em kW, tipos de plugue). Um eletroposto ultrarrápido (150 kW) exerce uma força de atração gravitacional e uma penalidade de concorrência significativamente maiores do que um carregador de parede (7 kW).
    \item \textbf{Modelos de escolha discreta (\textit{Logit/Probit}):} Substituir o decaimento linear da Zona de Transição por modelos probabilísticos de comportamento do consumidor, estimando a probabilidade real de um usuário desviar sua rota para um posto específico.
\end{itemize}

\subsection{Expansão do modelo de otimização (MILP)}

A estrutura puramente linear do modelo atual (\ref{sec:modelo}) permite a adição de novas restrições sem comprometer a viabilidade computacional. As principais extensões previstas incluem:

\begin{itemize}
    \item \textbf{Restrição de orçamento global (CAPEX):} Implementar uma inequação que limite o número de nós selecionados ($x_c$) ao orçamento de capital disponível, forçando o solver a priorizar os candidatos de mais alta eficiência marginal ($ \sum c_c \cdot x_c \leq B $).
    \item \textbf{Capacitação da rede elétrica:} Integrar restrições físicas da rede de distribuição. O modelo deve evitar a alocação de múltiplos eletropostos de alta potência em um mesmo alimentador local que não possua capacidade de carga ociosa suficiente, prevenindo instabilidades na rede.
    \item \textbf{Formulação multi-período:} Expandir o vetor de decisão para contemplar o horizonte temporal de adoção de VEs ($x_{c,t}$). O modelo não apenas decidirá \textit{onde} instalar, mas \textit{quando} instalar cada estação ao longo de uma década, acompanhando a curva projetada de demanda.
\end{itemize}

\subsection{Otimização robusta e incerteza}

Os dados de Ponto de Interesse (POI) representam uma fotografia estática da demanda urbana. No entanto, o tecido urbano é dinâmico. Trabalhos futuros poderão aplicar técnicas de \textbf{Otimização Robusta} para lidar com a incerteza nos parâmetros de atratividade ($A_c$). O objetivo será encontrar uma rede de eletropostos que mantenha um desempenho satisfatório e garanta a meta de cobertura ($Q_{\min}^{\text{cob}}$) mesmo sob os piores cenários de projeção de demanda ou falhas parciais na rede concorrente.

\vspace{0.5cm}
O desenvolvimento iterativo destas extensões não apenas aumentará a fidelidade do modelo frente à realidade operacional das cidades inteligentes, mas também consolidará a contribuição científica desta tese para a área de Pesquisa Operacional aplicada à mobilidade elétrica.

\end{document}